% --- Archon physics formalization source begin ---
% archon:physics
% archon:covers PhyXMiniProblems/problem_phyx_mini_0168.lean
% archon:source-report reports/phyx_mini/problem_phyx_mini_0168.source.json

\chapter{Physics problem phyx\_mini\_0168}
\label{ch:PhyXMiniProblems_problem_phyx_mini_0168}

\paragraph{Problem source.}
\#\# Physical scenario

A submerged scuba diver hears the sound of a boat horn directly above her on the surface of the lake. At the same time, a friend on dry land \textbackslash{}( 22.0\textbackslash{},\textbackslash{}text\{m\} \textbackslash{}) from the boat also hears the horn (\textbackslash{}textbf\{figure\}). The horn is \textbackslash{}( 1.2\textbackslash{},\textbackslash{}text\{m\} \textbackslash{}) above the surface of the water. Both air and water are at \textbackslash{}( 20\textasciicircum{}\{\textbackslash{}circ\}\textbackslash{}text\{C\} \textbackslash{}).

\#\# Question

What is the distance (labeled “?”) from the horn to the diver?

\#\# Answer choices

- A: A: 90.8m

- B: B: 94.8m

- C: C: 85.8m

- D: D: 65.8m

\#\# Figure caption (auxiliary; use the image as primary evidence)

The image depicts an underwater scene involving a diver, a boat, and a person on the shore. 

- A diver is underwater, with an arrow pointing down to them marked with a question mark, indicating an unknown depth.
- A figure stands on the shore, and a horizontal line from their feet extends to the edge of the water. The distance from this person to the water's edge is labeled "22.0 m."
- To the right, part of a boat is visible above the water. A vertical line with an arrow extends downward from the boat into the water, indicating the position directly above the diver.
- The vertical line and arrow are aligned with the boat's position on the water, suggesting a visual or auditory path relevant to the diver's position.
- The illustration contains elements that suggest the situation might relate to measuring depths or distances in a marine context.

\#\# Dataset metadata

Sound; Physical Model Grounding Reasoning, Spatial Relation Reasoning

\paragraph{Recorded answer/context.}
A: A: 90.8m

\paragraph{Figure/image path.}
/root/proposal\_for\_physic/science-mango/phyx\_mini\_run/phyx\_data/test\_image/168.png

\paragraph{Formalization target.}
create a compiling Lean file with sorry bodies at `PhyXMiniProblems/problem\_phyx\_mini\_0168.lean`. The Lean declarations must preserve the physical quantities, dimensions or dimensional roles, figure labels, governing-law hypotheses, and final relation expressed by this problem.
Use Mathlib/Physlib names found through LeanExplore where available. If a physics API is missing, introduce faithful local abstractions rather than scalar placeholder aliases.

\begin{theorem}[Physics formalization target]
\label{thm:physics:phyx_mini_0168:target}
\lean{PhyXMiniProblems.ProblemPhyXMini0168.hornToDiverDistance_rounds_to_answerA}
\uses{def:physics:phyx-mini-0168:phyxminiproblems-problemphyxmini0168-metersvalue, def:physics:phyx-mini-0168:phyxminiproblems-problemphyxmini0168-boathorndiversetup, def:physics:phyx-mini-0168:phyxminiproblems-problemphyxmini0168-matchesproblemandfigure, def:physics:phyx-mini-0168:phyxminiproblems-problemphyxmini0168-matchessoundspeeddataat20c, def:physics:phyx-mini-0168:phyxminiproblems-problemphyxmini0168-hasphysicalparameters, def:physics:phyx-mini-0168:phyxminiproblems-problemphyxmini0168-satisfiessoundpropagationlaws, def:physics:phyx-mini-0168:phyxminiproblems-problemphyxmini0168-isclosestanswerchoice}
The assigned autoformalize agent should translate this physics problem into Lean declarations in the covered file, with theorem and lemma proof bodies written as `by sorry`.
\end{theorem}
\begin{proof}
This is an autoformalization task, not a proof task. Produce statements that can later be proved by the physics prover without weakening the physical model.
\end{proof}

% --- Archon declaration topology begin ---
\section{Declaration topology}
\begin{definition}[Acoustic Length]
  \label{def:physics:phyx-mini-0168:phyxminiproblems-problemphyxmini0168-acousticlength}
  \lean{PhyXMiniProblems.ProblemPhyXMini0168.AcousticLength}
  A physical acoustic-path length, independent of the chosen unit readout
\end{definition}
\begin{proof}
  This is the declared model definition of Acoustic Length.
\end{proof}
\begin{definition}[Acoustic Duration]
  \label{def:physics:phyx-mini-0168:phyxminiproblems-problemphyxmini0168-acousticduration}
  \lean{PhyXMiniProblems.ProblemPhyXMini0168.AcousticDuration}
  A physical propagation duration, independent of the chosen unit readout
\end{definition}
\begin{proof}
  This is the declared model definition of Acoustic Duration.
\end{proof}
\begin{definition}[Acoustic Speed]
  \label{def:physics:phyx-mini-0168:phyxminiproblems-problemphyxmini0168-acousticspeed}
  \lean{PhyXMiniProblems.ProblemPhyXMini0168.AcousticSpeed}
  A physical sound speed, carrying the length-per-time dimension
\end{definition}
\begin{proof}
  This is the declared model definition of Acoustic Speed.
\end{proof}
\begin{definition}[meters Value]
  \label{def:physics:phyx-mini-0168:phyxminiproblems-problemphyxmini0168-metersvalue}
  \lean{PhyXMiniProblems.ProblemPhyXMini0168.metersValue}
  \uses{def:physics:phyx-mini-0168:phyxminiproblems-problemphyxmini0168-acousticlength}
  Read a physical length as a real number of metres
\end{definition}
\begin{proof}
  This is the declared model definition of meters Value.
\end{proof}
\begin{definition}[seconds Value]
  \label{def:physics:phyx-mini-0168:phyxminiproblems-problemphyxmini0168-secondsvalue}
  \lean{PhyXMiniProblems.ProblemPhyXMini0168.secondsValue}
  \uses{def:physics:phyx-mini-0168:phyxminiproblems-problemphyxmini0168-acousticduration}
  Read a physical duration as a real number of seconds
\end{definition}
\begin{proof}
  This is the declared model definition of seconds Value.
\end{proof}
\begin{definition}[meters Per Second Value]
  \label{def:physics:phyx-mini-0168:phyxminiproblems-problemphyxmini0168-meterspersecondvalue}
  \lean{PhyXMiniProblems.ProblemPhyXMini0168.metersPerSecondValue}
  \uses{def:physics:phyx-mini-0168:phyxminiproblems-problemphyxmini0168-acousticspeed}
  Read a physical speed as a real number of metres per second
\end{definition}
\begin{proof}
  This is the declared model definition of meters Per Second Value.
\end{proof}
\begin{definition}[Sound Medium]
  \label{def:physics:phyx-mini-0168:phyxminiproblems-problemphyxmini0168-soundmedium}
  \lean{PhyXMiniProblems.ProblemPhyXMini0168.SoundMedium}
  The two media traversed by the sound in the problem
\end{definition}
\begin{proof}
  This is the declared model definition of Sound Medium.
\end{proof}
\begin{definition}[Sound Path Leg]
  \label{def:physics:phyx-mini-0168:phyxminiproblems-problemphyxmini0168-soundpathleg}
  \lean{PhyXMiniProblems.ProblemPhyXMini0168.SoundPathLeg}
  The three constant-medium path legs used to compare the two listeners' travel times. The latter two legs concatenate to form the vertical horn-to-diver path marked ? in the figure
\end{definition}
\begin{proof}
  This is the declared model definition of Sound Path Leg.
\end{proof}
\begin{definition}[Boat Horn Diver Setup]
  \label{def:physics:phyx-mini-0168:phyxminiproblems-problemphyxmini0168-boathorndiversetup}
  \lean{PhyXMiniProblems.ProblemPhyXMini0168.BoatHornDiverSetup}
  \uses{def:physics:phyx-mini-0168:phyxminiproblems-problemphyxmini0168-acousticlength, def:physics:phyx-mini-0168:phyxminiproblems-problemphyxmini0168-acousticduration, def:physics:phyx-mini-0168:phyxminiproblems-problemphyxmini0168-acousticspeed, def:physics:phyx-mini-0168:phyxminiproblems-problemphyxmini0168-soundmedium, def:physics:phyx-mini-0168:phyxminiproblems-problemphyxmini0168-soundpathleg}
  Physical quantities for the boat-horn experiment. The temperature function stores Celsius readouts, while lengths, times, and speeds themselves remain dimensionful quantities. In particular, no numerical value for the requested hornToDiverDistance is stored here
\end{definition}
\begin{proof}
  This is the declared model definition of Boat Horn Diver Setup.
\end{proof}
\begin{definition}[Matches Problem And Figure]
  \label{def:physics:phyx-mini-0168:phyxminiproblems-problemphyxmini0168-matchesproblemandfigure}
  \lean{PhyXMiniProblems.ProblemPhyXMini0168.MatchesProblemAndFigure}
  \uses{def:physics:phyx-mini-0168:phyxminiproblems-problemphyxmini0168-metersvalue, def:physics:phyx-mini-0168:phyxminiproblems-problemphyxmini0168-boathorndiversetup}
  Problem-text and primary-figure information. The direct horn-to-friend sound path is labeled 22.0 m; the horn is 1.2 m above the lake surface; and the diver is directly below the horn, so the requested vertical distance is the sum of its air and water legs. This does not state the requested numerical distance
\end{definition}
\begin{proof}
  This is the declared model definition of Matches Problem And Figure.
\end{proof}
\begin{definition}[Matches Sound Speed Data At20 C]
  \label{def:physics:phyx-mini-0168:phyxminiproblems-problemphyxmini0168-matchessoundspeeddataat20c}
  \lean{PhyXMiniProblems.ProblemPhyXMini0168.MatchesSoundSpeedDataAt20C}
  \uses{def:physics:phyx-mini-0168:phyxminiproblems-problemphyxmini0168-meterspersecondvalue, def:physics:phyx-mini-0168:phyxminiproblems-problemphyxmini0168-boathorndiversetup}
  Temperature and standard sound-speed readouts used for air and fresh water at 20 °C: 344 m/s in air and 1482 m/s in water. These are calibration data, not the requested distance
\end{definition}
\begin{proof}
  This is the declared model definition of Matches Sound Speed Data At20 C.
\end{proof}
\begin{definition}[Has Physical Parameters]
  \label{def:physics:phyx-mini-0168:phyxminiproblems-problemphyxmini0168-hasphysicalparameters}
  \lean{PhyXMiniProblems.ProblemPhyXMini0168.HasPhysicalParameters}
  \uses{def:physics:phyx-mini-0168:phyxminiproblems-problemphyxmini0168-metersvalue, def:physics:phyx-mini-0168:phyxminiproblems-problemphyxmini0168-secondsvalue, def:physics:phyx-mini-0168:phyxminiproblems-problemphyxmini0168-meterspersecondvalue, def:physics:phyx-mini-0168:phyxminiproblems-problemphyxmini0168-boathorndiversetup}
  Positivity and nonnegativity conditions for the physical setup
\end{definition}
\begin{proof}
  This is the declared model definition of Has Physical Parameters.
\end{proof}
\begin{definition}[Satisfies Sound Propagation Laws]
  \label{def:physics:phyx-mini-0168:phyxminiproblems-problemphyxmini0168-satisfiessoundpropagationlaws}
  \lean{PhyXMiniProblems.ProblemPhyXMini0168.SatisfiesSoundPropagationLaws}
  \uses{def:physics:phyx-mini-0168:phyxminiproblems-problemphyxmini0168-metersvalue, def:physics:phyx-mini-0168:phyxminiproblems-problemphyxmini0168-secondsvalue, def:physics:phyx-mini-0168:phyxminiproblems-problemphyxmini0168-meterspersecondvalue, def:physics:phyx-mini-0168:phyxminiproblems-problemphyxmini0168-boathorndiversetup}
  Governing propagation and reception-time laws. Constant-speed travel on each homogeneous leg gives distance = speed × time. Because the same horn emission reaches both listeners simultaneously, the air travel time to the friend equals the sum of the vertical air and water travel times to the diver. Neither law contains the numerical answer
\end{definition}
\begin{proof}
  This is the declared model definition of Satisfies Sound Propagation Laws.
\end{proof}
\begin{definition}[Answer Choice]
  \label{def:physics:phyx-mini-0168:phyxminiproblems-problemphyxmini0168-answerchoice}
  \lean{PhyXMiniProblems.ProblemPhyXMini0168.AnswerChoice}
  Labels of the four displayed multiple-choice answers
\end{definition}
\begin{proof}
  This is the declared model definition of Answer Choice.
\end{proof}
\begin{definition}[answer Choice Meters]
  \label{def:physics:phyx-mini-0168:phyxminiproblems-problemphyxmini0168-answerchoicemeters}
  \lean{PhyXMiniProblems.ProblemPhyXMini0168.answerChoiceMeters}
  \uses{def:physics:phyx-mini-0168:phyxminiproblems-problemphyxmini0168-answerchoice}
  Metre readout printed beside each displayed answer choice
\end{definition}
\begin{proof}
  This is the declared model definition of answer Choice Meters.
\end{proof}
\begin{definition}[Is Closest Answer Choice]
  \label{def:physics:phyx-mini-0168:phyxminiproblems-problemphyxmini0168-isclosestanswerchoice}
  \lean{PhyXMiniProblems.ProblemPhyXMini0168.IsClosestAnswerChoice}
  \uses{def:physics:phyx-mini-0168:phyxminiproblems-problemphyxmini0168-answerchoice, def:physics:phyx-mini-0168:phyxminiproblems-problemphyxmini0168-answerchoicemeters}
  A displayed choice is uniquely closest to an actual metre readout
\end{definition}
\begin{proof}
  This is the declared model definition of Is Closest Answer Choice.
\end{proof}
\begin{lemma}[horn To Diver Distance exact]
  \label{lem:physics:phyx-mini-0168:phyxminiproblems-problemphyxmini0168-horntodiverdistance-exact}
  \lean{PhyXMiniProblems.ProblemPhyXMini0168.hornToDiverDistance_exact}
  \uses{def:physics:phyx-mini-0168:phyxminiproblems-problemphyxmini0168-metersvalue, def:physics:phyx-mini-0168:phyxminiproblems-problemphyxmini0168-boathorndiversetup, def:physics:phyx-mini-0168:phyxminiproblems-problemphyxmini0168-matchesproblemandfigure, def:physics:phyx-mini-0168:phyxminiproblems-problemphyxmini0168-matchessoundspeeddataat20c, def:physics:phyx-mini-0168:phyxminiproblems-problemphyxmini0168-hasphysicalparameters, def:physics:phyx-mini-0168:phyxminiproblems-problemphyxmini0168-satisfiessoundpropagationlaws}
  Before rounding to the precision of the choices, the idealized equal-arrival model gives 19524 / 215 m for the full horn-to-diver distance
\end{lemma}
\begin{proof}
  The conclusion follows from the governing relations and figure readouts named in the statement of horn To Diver Distance exact.
\end{proof}
% --- Archon declaration topology end ---
% --- Archon physics formalization source end ---
